\section{Social Graph \texorpdfstring{\emojisoctitle}{Social Graph}}
\label{sec-soc}

The social graph typically consists of entities connected by their social relations and is ubiquitous across real-world applications. A primary example is social networks like Twitter and Facebook, where entities represent individuals linked by social interactions (e.g., friendships, followers/followees, likes, and mentions). These social graphs extend beyond human interactions and are not confined to living entities, such as tortoises co-using the same burrow in animal social networks~\cite{rossi2015network}, complementary products co-purchased by the same customer in E-commerce recommender systems~\cite{wang2024knowledge2, wu2024stark}, or even the LLM-simulated social agents~\cite{zhang2023exploring, li2023metaagents}. The wealth of social-relational knowledge in these social graphs is the golden resource for GraphRAG\cite{zeng2024large, jiang2023social, xie2023factual, wang2024augmenting, zeng2024federated, wang2024knowledge2, deldjoo2024review, du2024large, wu2024stark, huang2021graph, zhang2021learning, qiu2020exploiting, wei2024llmrec, kim2020retrieval, tian2021joint}
, which is reviewed in this section.


\subsection{Application Tasks}

\begin{itemize}[leftmargin=*]
    \item \textbf{Entity Property Prediction}: 
    Entity property prediction focuses on predicting properties and classifying categories for social entities in social networks, examples of which include the prediction of partnership compatibility, assessment of morality, detection of account suspensions, identification of toxic behaviors~\cite{jiang2023social} and product property prediction~\cite{wang2024knowledge2}.


    \item \textbf{Text Generation}: 
    Text generation aims to produce text that aligns with social contexts and norms. Typically, the interplay between structural and textual information, such as proximity-based network homophily and role-based similarity~\cite{ahmed2018learning}, serves as the foundation for text generation. For instance, \citet{wang2024augmenting} retrieves the texts of proximity/role-similar nodes to enhance the text generation of the target node. \citet{kim2020retrieval, xie2023factual} generate personalized recommendation explanations by retrieving customers'/products' historical reviews. In addition, some other works also leverage semantic similarity to retrieve reference/attributes/opinions and augment the downstream review generation~\cite{sharma2018cyclegen, dong2017learning, park2015retrieval}.
    
    

    \item \textbf{Recommendation}:
    The recommendation task aims to find the most relevant items to satisfy user demands. Due to the missing and sparse customer/item interactions (e.g., cold-start issues), GraphRAG can be naturally applied to boost customer/item sparse interaction by retrieving additional meta-knowledge. Depending on the concrete recommendation scenario, GraphRAG has been used for graph-based recommendation~\cite{wei2024llmrec, du2024large}, next-item recommendation~\cite{wang2024knowledge2, huang2021graph, zhang2021learning, qiu2020exploiting, zeng2024federated, wang2023zero}, and conversational recommendation~\cite{friedman2023leveraging}
    
    
    
    \item \textbf{Question-answering}: Question-answering tasks are encountered not only in knowledge and document graph domains but also in social graphs. For example, a user might ask, "What are the best parks for family gatherings around Los Gatos?" and expect a personalized query answer~\cite{zeng2024large, kemper2024retrieval}. Furthermore, such questions might explicitly require graph-structured reasoning. For instance, \citet{wu2024stark} build a semi-structured knowledge base, and some queries might ask, "Can you list the products made by Nike?" Answering this question requires not only a deep understanding of the query but also familiarity with the structural information of the data.


    \item \textbf{Fake News Detection}: Detecting fake news necessitates considering both semantic content (e.g., the content of the news) and structural interactions (e.g., interactions among news). For instance, \citet{ram2024credirag} evaluates the credibility of a Reddit post based on the credibility of other Reddit posts retrieved by their common interactions with the same common commenters.
    
\end{itemize}


\subsection{Social Graph Construction}
The relations that GraphRAG leverages to derive additional information from social graphs can be mainly summarized into three rationales: proximity-based, role-based, and personalization-based rationale. The proximity-based rationale, rooted in the adage "birds of a feather flock together," suggests that nodes close to each other in a social network often share similar properties~\cite{mcpherson2001birds}. For example, close friends tend to possess similar hobbies. The role-based rationale focuses on nodes with similar local subgraph structures sharing similar features or label distributions~\cite{donnat2018learning}. For example, managers at the same hierarchical level within a company generally have comparable job titles and responsibilities, and hub airports exhibit similar operational characteristics and strategic importance. Lastly, the personalization-based rationale refers to individuals' uniqueness regarding their characteristics and interactions. For example, in recommender systems, users interact with items in various ways, such as clicking, viewing, adding to a cart, purchasing, and reviewing, each interaction providing valuable knowledge that GraphRAG can leverage to personalize generated content. Based on the above three relational rationales in the social graphs, the social graph construction methods can be summarized as follows:

\begin{itemize}[leftmargin=*]
    \item \textbf{User-User-Interaction}~\cite{jiang2023social}: This type of social graph represents user-to-user interactions commonly seen on social networks such as Twitter, Reddit, and Facebook. Examples include follower-followee relationships on Twitter, friendship relationships on Facebook, and user comments on other users' threads on Reddit. Note that this user-user interaction naturally exists in the real world without human curation or modification compared to documents or knowledge graphs.

    \item \textbf{User-Item-Interaction}~\cite{wang2024knowledge2, wang2024augmenting, xie2023factual, kim2020retrieval}: This type of social graph represents user-to-item interactions commonly found on e-commerce platforms like Amazon and eBay. These interactions, including purchasing, adding-to-the-cart, and viewing, can be modeled as a bipartite graph, where each type of interaction reflects a unique user intention.

    \item \textbf{Item-Item-Interaction}~\cite{wang2024knowledge2, zhang2021learning, wu2024stark}: This social graph captures interactions between items, typically identified by shared interactions from the same customer or user. For example, a "co-view" interaction between two products indicates that they are viewed by the same customer, while a "view-add-to-cart" interaction indicates that one product is first viewed and the other is added to the cart next by the same customer. In e-commerce networks, these co-interactions between two items can be broadly categorized into complementary and substitute relationships~\cite{zhao2024recommender}.

    \item \textbf{Metadata-Interaction}~\cite{wu2024stark, yang2024common}: Items and users often possess metadata. For example, products on platforms like Amazon may include brand, manufacturer, and color attributes. This metadata can be represented as additional node types and the corresponding edges indicating relations between products and attributes, such as ownership or association.

    \item \textbf{Agent-Agent Interaction}~\cite{chen2023agentverse, zhang2023exploring}: With the increasing intelligence of LLMs, recent literature has explored the potential of using LLM-powered agents to simulate social behaviors, such as collaboration, debate, and reflection. These agent-agent interactions could also be used for GraphRAG.
    
\end{itemize}

Note that among the five types of interaction mentioned above, the user-user, user-item, and metadata relations are naturally formulated without human curation. In contrast, item-item interactions are generated by manual extraction, and agent-agent interactions require simulations.



\subsection{Retriever}
The relations of the social graphs constructed according to the above methods can be leveraged in GraphRAG to enhance downstream tasks by retrieving additional information. For example, retrieving historical metadata and customer interactions can improve recommendations. Next, we review representative retrieval methods in GraphRAG for social graphs.

\begin{itemize}[leftmargin=*]
    \item \textbf{ID-based Retriever}: Similar to entity linking, the ID-based retriever works by retrieving content specifically generated by a user/item, examples of which include retrieving historical item interactions of a specific customer~\cite{zeng2024large, jiang2023social, zeng2024federated}, reviews posted on a specific product~\cite{xie2023factual}, and meta-data information about a specific user~\cite{wu2024stark, salemi2023lamp}.

    \item \textbf{Filtering-based Retriever}: Based on the ID-based retriever, the Filtering-based Retriever retrieves additional content based on collaborative filtering~\cite{wang2023zero}. For user filtering, Top K users are identified by comparing the similarity of their historical item interactions with those of the target user demanding recommendation. To enhance the context of the target user, it retrieves the most popular items from those Top K users. In contrast, for item filtering, it identifies the Top K items that are most similar to the current item by comparing their user-interaction history. Among them, the most popular ones would be fetched to augment the current item recommendation.


    \item \textbf{Social Relational Retriever}: Like the ID-based Retriever, the Social Relational Retriever focuses on retrieving knowledge from entities sharing certain relations with the target entities on hand. For example, \citet{du2024large} hierarchically retrieve neighboring texts several hops from the central node to augment the semantic information of the target user/item. Meanwhile, \citet{wang2024augmenting} retrieve texts corresponding to nodes that share high proximity-based and role-based similarity, respectively.
    
    \item \textbf{Integrated Neural-Symbolic Retriever}: This approach leverages both symbolic and neural retrievers to improve retrieval effectiveness~\cite{wang2024knowledge2, zeng2024federated, wang2020global}. The symbolic retriever retrieves information by following explicitly defined rules, such as retrieving based on identifiers, structured relationships, or interaction patterns, ensuring that the retrieved data strictly aligns with specific criteria. Meanwhile, the neural retriever complements this by using embedding-based similarity, capturing nuanced patterns and contextual relationships that may not be directly encoded in rules. Integrating them together provides a better trade-off between rule-based precision and neural-based adaptability and generalizability. For instance, \citet{wang2020global, wang2024knowledge2, huang2021graph, qiu2020exploiting} first retrieve K-hop neighboring products from the product knowledge graph (symbolic retriever) for products in the user session and further utilize neural-based adaptive filtering to aggregate items that are most relevant to the current sequence (i.e., neural retriever). Similarly, \citet{zeng2024federated} address data sparsity and heterogeneity by combining ID-based retrieval that retrieves movies based on the user ID of one party with a text-based retriever that enables movie retrieval between parties.

\end{itemize}


\subsection{Organizer}
To further enhance the retrieved content, the organizer for the social graphs employs specialized techniques beyond the typical re-ranking and filtering used in other graph domains~\cite{zeng2024large, hou2024large, deldjoo2024review}. For social graphs, the organizer of GraphRAG often uses Keyword Extraction, Profile Summarization, and Hierarchical Graph Aggregation to curate the retrieved content:

\begin{itemize}[leftmargin=*]
    \item \textbf{Keyword Extraction}: Keyword extraction identifies the most relevant and informative keywords from the retrieved content. These extracted keywords guide the downstream generator in prioritizing attention and reduce the risk of overwhelming LLMs with excessive context. For example, \citet{xie2023factual} use an embedding estimator to pinpoint keywords aligned with a personalized latent query, which are then used to generate explanations.

    \item \textbf{Profile Summarization}: Profile summarization creates rich and detailed user profiles that capture attributes such as age, gender, preferred and disliked genres, favorite directors, country, and language derived from users' past interactions and item metadata~\cite{wei2024llmrec}. This enriched user data increases the informativeness of the retrieved content while safeguarding privacy by using rewritten profiles. Additionally, after retrieving related movies via user/item filtering, \citet{wang2023zero} apply a three-step prompting strategy to extract features tailored to the user and select representative movies by directly instructing LLMs. These extracted features and representative movies are further used as the user profile to guide the final recommendation generation. \citet{guo2024lightrag} leverage LLMs to generate more details about the entity profile based on its keywords/short phrases from external data to aid text generation.

    \item \textbf{Hierarchical Graph Aggregation and Summarization}: 
    When retrieving neighborhood textual information, the exponentially expanding receptive field at higher neighborhood layers often causes the aggregated content to exceed manageable limits, potentially diminishing the LLM's effectiveness~\cite{li2024long}. This challenge highlights the need for hierarchical graph aggregation and summarization. The primary idea is to first summarize neighborhood information retrieved at each layer before propagating it to the next layer. This approach keeps the volume of text each node receives within consistent bounds. For instance, \citet{du2024large} recursively aggregate information from neighboring nodes at higher layers and leverage LLMs to rephrase and compress the content before sharing it with relevant nodes in lower layers. Consequently, this strategy expands the receptive field while optimizing the computational budget for downstream generation tasks.
\end{itemize}
 
\subsection{Generator}
Once the relevant content is retrieved and organized appropriately, it is further processed by a downstream generator to produce the final content. Depending on the desired output, existing generators used for social graphs can be categorized into LLM-based text generators and Prediction-based generators. The choice between these two depends on the format of the desired output and the requirements of the social graph applications.

\begin{itemize}[leftmargin=*]
    \item \textbf{LLM-based Text Generation}: This generator is typically used when the downstream application requires text outputs~\cite{zeng2024federated, xie2023factual, wang2023zero, wang2024augmenting}, such as item recommendations based on names, recommendation explanations, and review generation. Due to the probabilistic nature of the next token prediction, the generated text may not always precisely align with the desired output. To address this hallucination issue, the ground-truth text is often used for grounding the generated texts, such as matching generated items with the ground-truth items on the platform~\cite{hou2024large}.

    \item \textbf{Prediction-based Generation}: The generator directly predicts outputs and is primarily used for non-textual tasks~\cite{du2024large}, such as item recommendation and social prediction. For instance, Graph Neural Networks have become popular choices for graph-based recommendation~\cite{du2024large, wei2024llmrec} and transformers are used for item recommendation tasks~\cite{wang2024knowledge2}. Furthermore, in social property prediction~\cite{jiang2023social}, the generator could simply be a multilayer perception used for classification (e.g., partisanship classification) or regression (e.g., morality regression).
    
\end{itemize}


\subsection{Resources and Tools}

This section lists common resources and tools used for GraphRAG on social graphs. The summary and the link of  are itemized as follows:

\subsubsection{Data Resources}
\begin{itemize}[leftmargin=*]
    \item \textbf{STARK-Amazon}\footnote{\href{https://stark.stanford.edu/dataset_amazon.html}{https://stark.stanford.edu/dataset\_amazon.html}}~\cite{wu2024stark}: STARK-Amazon is a large-scale, semi-structured retrieval benchmark dataset for product search on the Amazon platform, integrating textual and relational knowledge bases. The nodes in the dataset represent products, colors, brands, and categories, while edges capture relationships like also\_bought, also\_viewed, has\_brand, has\_category, and has\_color. Rich textual information, including product descriptions, customer reviews, and entity names, provides valuable context for retrieval tasks. Queries are generated by sampling relational templates and grounding them with specific entities, then leveraging LLMs to synthesize relevant textual and relational information. This process results in queries that capture customer interests, interpret specialized descriptions, and deduce relationships involving multiple entities within the query.
    
    \item \textbf{Amazon-Review}\footnote{\href{https://jmcauley.ucsd.edu/data/amazon/}{https://jmcauley.ucsd.edu/data/amazon/}}~\cite{ni2019justifying, he2016ups, mcauley2015image, xie2023factual, zeng2024federated}: The Amazon-Review dataset includes reviews with ratings, text, and helpfulness votes, along with product metadata such as descriptions, categories, price, brand, and image features. Additionally, it provides link information through “also viewed” and “also bought” graphs. This dataset has two versions: the initial release logging the review data between 1996 and 2014~\cite{he2016ups, mcauley2015image}, and a more recent version that logs ongoing data in 2014~\cite{ni2019justifying}. The latest version also adds new metadata, including detailed product information, bullet points, and an ethical details table. Extensive GraphRAG work~\cite{} in social graphs has leveraged this dataset to build RAG frameworks for recommendation research.
    
    \item \textbf{MovieLens}\footnote{\href{https://grouplens.org/datasets/movielens/}{https://grouplens.org/datasets/movielens/}}~\cite{wang2023zero, wei2024llmrec, zeng2024federated}: The MovieLens datasets describe people’s expressed preferences for movies~\cite{harper2015movielens}. These preferences take the form of tuples, each resulting in a person expressing a preference (a 0-5 star rating) for a movie at a particular time. These preferences were entered through the MovieLens website — a recommender system that asks its users to give movie ratings to receive personalized movie recommendations. Side information includes movie title, year, and genre in textual format. \cite{wei2024llmrec} further crawl the visual content of movie posters accessible from \href{https://github.com/HKUDS/LLMRec}{here}. Note that Movielen-100k is frequently used~\cite{wei2024llmrec, zeng2024federated} and provides users with demographic information such as user ID, age, gender, occupation, and zip code.

    \item \textbf{Netflix}\footnote{\href{https://www.kaggle.com/datasets/netflix-inc/netflix-prize-data}{https://www.kaggle.com/datasets/netflix-inc/netflix-prize-data}}~\cite{wei2024llmrec}: This dataset was constructed to support participants in the Netflix Prize. The movie rating files contain over 100 million ratings from 480 thousand randomly chosen, anonymous Netflix customers over 17 thousand movie titles.  The data were collected between October 1998 and December 2005 and reflect the distribution of all ratings received during this period. The date of each rating, the title, and the year of release for each movie ID are also provided. The multi-model side information is collected through web crawling~\cite{wei2024llmrec}, which is further stored \href{https://github.com/HKUDS/LLMRec}{here}.

    \item \textbf{Yelp}\footnote{\href{https://www.yelp.com/dataset}{https://www.yelp.com/dataset}}~\cite{xie2023factual}: The Yelp dataset is a rich resource for academic research, particularly in fields like data science, natural language processing (NLP), and machine learning. This comprehensive dataset includes over 6.9 million reviews, 150,346 businesses, and 200,100 photos, covering 11 metropolitan areas. Researchers and students can explore real-world business and user data, with extensive details about businesses such as location, categories, attributes (like hours, parking, and ambiance), and aggregated check-ins. Reviews include full-text content, user ratings, and metadata, while user profiles offer insights into social interactions and behaviors, including friends, compliments, and review history. Additionally, the dataset contains 908,915 tips, providing quick recommendations and insights, and 1.2 million business attributes, offering granular information about service offerings.


    \item \textbf{Weibo}\footnote{\href{https://www.aminer.cn/influencelocality}{https://www.aminer.cn/influencelocality}}~\cite{zhang2013social}: The Weibo dataset offers a rich snapshot of user interactions, behaviors and social connections within the Sina Weibo platform, capturing both static and dynamic facets of the social network. Starting with 100 randomly selected seed users, the dataset scales to encompass 1.7 million users and around 0.4 billion following relationships, averaging 200 followers per user. Each user is characterized by social profile details, including name, gender, verification status, and follower/followee counts. Tweet content is provided in both original Chinese and indexed formats, supporting robust research in social network analysis, content diffusion, and user interaction dynamics, making it an invaluable resource for RAG studies.


    \item \textbf{Brexit}\footnote{\href{https://github.com/somethingx01/TopicalAttentionBrexit?tab=readme-ov-file}{https://github.com/somethingx01/TopicalAttentionBrexit?tab=readme-ov-file}}~\cite{zhu2020neural}: This dataset includes a portion of the X (Twitter) network, specifically the remain-leave discourse before the 2016 UK Referendum on exiting the EU. It comprises a network with 7,589 users, 532,459 directed follow relationships, and 19,963 tweets, each associated with a binary stance. The dataset is preprocessed according to \cite{minici2022cascade} to assign each user a scalar value between 0 and 1, referred to as opinion, representing the average stance of the tweets retweeted by the user. The stance of each tweet is either 0 (“Remain”) or 1 (“Leave”).


    \item \textbf{Diginetica}\footnote{\href{https://competitions.codalab.org/competitions/11161\#learn_the_details-data2}{https://competitions.codalab.org/competitions/11161\#learn\_the\_details-data2}}~\cite{wang2024knowledge2}: This dataset comprises user session logs from an e-commerce search engine. The data spans six months and captures user interactions, including clicks, product views, and purchases. Each user session, defined by a one-hour inactivity period, contains anonymized user IDs, hashed queries, product descriptions, metadata (price, hashed product names, image identifiers, and product categories), and log-scaled prices.

    \item \textbf{Yoochoose}\footnote{\href{https://www.kaggle.com/datasets/chadgostopp/recsys-challenge-2015}{https://www.kaggle.com/datasets/chadgostopp/recsys-challenge-2015}}~\cite{wang2024knowledge2}: The Yoochoose dataset contains session data from an online European retailer. Each session records the user's click events, with some sessions also including purchase events. Collected over several months in 2014, the data captures user interactions with the retailer's website. The product meta-data includes categories.
\end{itemize}

\subsubsection{Tools} 
\begin{itemize}[leftmargin=*]
    \item \textbf{X-Developer Platform}\footnote{\href{https://developer.x.com/en/docs/x-api/getting-started/about-x-api}{https://developer.x.com/en/docs/x-api/getting-started/about-x-api}}: The X Developer Platform provides powerful tools and resources for developers to integrate X's real-time, historical, and global data into their own applications. With three main products—X API, X Ads API, and X for Websites—the platform supports a wide range of use cases, from retrieving and analyzing Tweets to managing ad campaigns and embedding X content directly into websites. The X API offers endpoints for managing conversations, exploring trends, and engaging with users, while the X Ads API enables businesses to manage ads with custom targeting and analytics. X for Websites allows seamless embedding of live content to enhance website engagement. Through comprehensive documentation, libraries, and community support, the X Developer Platform enables developers to create innovative solutions using X’s data and engagement tools.
    
    \item \textbf{Reddit-API}\footnote{\href{https://support.reddithelp.com/hc/en-us/articles/16160319875092-Reddit-Data-API-Wiki}{https://support.reddithelp.com/hc/en-us/articles/16160319875092-Reddit-Data-API-Wiki}}: Reddit is a news aggregation and discussion platform where posts are organized into "subreddits," user-created boards moderated by the community. The Reddit API provides developers access to the site’s extensive collection of posts and comments. This free API has enabled the development of moderation tools, third-party applications, and training datasets for LLMs such as ChatGPT, Google, and Gemini. By using this API, we can query tremendous user interactions (i.e., comments and posts) and their corresponding textual contents (i.e., subreddit threads), which serve as golden resources for GraphRAG.
    
    \item \textbf{Rec-Bole}\footnote{\href{https://recbole.io/}{https://recbole.io/}}~\cite{zhao2021recbole}: RecBole is an open-source, unified, and comprehensive library designed for developing and benchmarking recommendation algorithms. Built with Python and PyTorch, RecBole offers researchers a streamlined, efficient framework to experiment with over 100 recommendation algorithms across four main types: General, Sequential, Context-Aware, and Knowledge-Based Recommendations. The platform simplifies data handling by providing pre-processed copies of 43 benchmark datasets, making it easy for users to dive into model testing and development. The provided user-product interactions, as well as text-formatted user/product meta-data, could also be used for GraphRAG.
\end{itemize}
